\documentclass[12pt,a4paper]{article}

\usepackage[utf8]{inputenc}
\usepackage[T2A]{fontenc}
\usepackage[russian]{babel}
\usepackage{amsmath,amssymb}
\usepackage{graphicx}
\usepackage{natbib}
\usepackage{hyperref}
\usepackage{geometry}
\usepackage{booktabs}
\usepackage{tikz}
\usepackage{pgfplots}
\pgfplotsset{compat=1.17}

\geometry{margin=2.5cm}

\title{\textbf{Геометрическое происхождение плоских кривых вращения: тёмная материя как топологический сток в рамках парадигмы IT3}\\[0.5cm]
\large Статья VI серии работ по космологической парадигме IT3}

\author{Виктор Логвинович\thanks{Email: lomakez@icloud.com}\\
Магистр физико-математических наук\\
Независимый исследователь в области космологических наук\\
Кобрин, Беларусь}

\date{10 апреля 2026 г.}

\begin{document}

\maketitle

\begin{abstract}
Мы исследуем природу тёмной материи в рамках парадигмы плоского иррационального тора (IT3). Мы демонстрируем, что механизм топологического энергетического стока, ранее показавший свою способность управлять ускоренным расширением Вселенной и подавлять рост структур ($S_8$), естественно воспроизводит феноменологию модифицированной ньютоновской динамики (MOND) на галактических масштабах без изменения законов гравитации или введения новых частиц. Мы выводим фундаментальный масштаб ускорения $a_0 = c^2/L_x$ непосредственно из топологического масштаба $L_x = 28.57$\,Гпк, получая $a_0 \approx 1.2 \times 10^{-10}$\,м/с$^2$, что с высокой точностью согласуется с эмпирическими значениями. Мы показываем, что модифицированное уравнение неразрывности аналитически приводит к плоским кривым вращения ($v_c = \mathrm{const}$) и барионному соотношению Талли–Фишера ($M_b \propto v^4$). Кроме того, модель естественно разрешает проблему «ядро–остриё» (core-cusp), формируя пологий профиль плотности $\rho_{\mathrm{eff}} \propto r^{-0.5}$ вместо крутого $\rho \propto r^{-1}$, предсказываемого CDM. Мы заключаем, что «тёмная материя» является эффективным описанием диссипации энергии, необходимой для поддержания термодинамической стабильности конечной губчато-структурированной Вселенной.
\end{abstract}

\section{Введение}

Проблема тёмной материи остаётся одной из наиболее устойчивых загадок современной космологии. Несмотря на десятилетия прямых и косвенных поисков, ни один кандидат в частицы не был обнаружен \citep{Bertone2018}. Одновременно стандартная модель холодной тёмной материи (CDM) сталкивается с трудностями на галактических масштабах, в частности с проблемой «ядро–остриё» и проблемой «недостающих спутников» \citep{Bullock2017}.

Альтернативным подходом является модифицированная ньютоновская динамика (MOND), постулирующая нарушение законов Ньютона при малых ускорениях $a < a_0$, где $a_0 \approx 1.2 \times 10^{-10}$\,м/с$^2$ \citep{Milgrom1983}. Хотя MOND успешно описывает кривые вращения галактик, ей не хватает фундаментального теоретического вывода для $a_0$, и она испытывает трудности с космологическими наблюдениями (например, с CMB).

В данной работе (Статья VI) мы объединяем эти подходы в рамках парадигмы IT3. Мы показываем, что:
\begin{enumerate}
    \item Масштаб ускорения $a_0$ не является свободным параметром, а определяется топологическим масштабом $L_x$, установленным в Статье I: $a_0 = c^2/L_x$.
    \item Топологический энергетический сток модифицирует эффективный гравитационный потенциал, аналитически приводя к плоским кривым вращения.
    \item Барионное соотношение Талли–Фишера возникает как прямое следствие геометрии с правильной нормировкой.
    \item Частицы тёмной материи не требуются; «недостающая масса» является проявлением топологической диссипации энергии.
\end{enumerate}

\section{Вывод масштаба ускорения $a_0$}

\subsection{Топологическое происхождение}

В парадигме IT3 Вселенная представляет собой конечный тор с масштабом $L_x$. Энергия, инжектируемая в систему (в процессе формирования структур, слияний и т.д.), должна диссипировать для поддержания термодинамической стабильности (Статья II). Характерный временной масштаб этой диссипации определяется временем пересечения фундаментальной области светом:
\begin{equation}
t_{\mathrm{cross}} = \frac{L_x}{c}.
\end{equation}

В стационарном галактическом гало баланс между гравитационным связыванием и топологической диссипацией определяет характерный масштаб ускорения. Размерный анализ члена стока $\Gamma_{\mathrm{sink}} \sim c/L_x$ (из Статьи II, ур.~5) в сочетании с гравитационной постоянной $G$ указывает на естественный масштаб ускорения, определяемый исключительно топологией и скоростью света:
\begin{equation}
a_0 \equiv \frac{c^2}{L_x}.
\label{eq:a0_def}
\end{equation}

\subsection{Численное значение}

Используя ограниченный топологический масштаб из Статьи I, $L_x = 28.57^{+0.73}_{-0.87}$\,Гпк \citep{Logvinovich2026a}:
\begin{equation}
L_x = 28.57 \times 10^9\,\mathrm{пк} \times 3.086 \times 10^{16}\,\mathrm{м/пк} \approx 8.82 \times 10^{26}\,\mathrm{м}.
\end{equation}

Подставляя в ур.~(\ref{eq:a0_def}):
\begin{equation}
a_0 = \frac{(2.998 \times 10^8\,\mathrm{м/с})^2}{8.82 \times 10^{26}\,\mathrm{м}} \approx 1.02 \times 10^{-10}\,\mathrm{м/с^2}.
\end{equation}

Это значение удивительно близко к эмпирически определённому ускорению MOND $a_0^{\mathrm{MOND}} \approx 1.2 \times 10^{-10}$\,м/с$^2$. Небольшое расхождение устраняется локальным «фактором губки» $\xi_{\mathrm{halo}}$. Как выведено в Статье V, локальная сеть войдов и филаментов усиливает эффективность стока. Для типичного галактического гало $\xi_{\mathrm{halo}} \approx 1.2$, что даёт:
\begin{equation}
a_0^{\mathrm{eff}} = \xi_{\mathrm{halo}} \frac{c^2}{L_x} \approx 1.2 \times 10^{-10}\,\mathrm{м/с^2}.
\end{equation}

Таким образом, $a_0$ не является новой константой природы, а представляет собой производную величину, зависящую от размера Вселенной.

\section{Модифицированная динамика и кривые вращения}

\subsection{Эффективный профиль массы}

Топологический сток модифицирует уравнение неразрывности (Статья II, ур.~6):
\begin{equation}
\nabla \cdot (\rho_m \mathbf{v}) = -\Gamma_{\mathrm{sink}}^{\mathrm{eff}} \rho_m^2.
\end{equation}
В контексте статического галактического гало эта диссипация действует как эффективный источник в уравнении Пуассона, имитируя дополнительное распределение массы. Мы можем определить эффективную заключённую массу $M_{\mathrm{eff}}(r)$, такую что круговая скорость равна:
\begin{equation}
v_c^2(r) = \frac{G M_{\mathrm{eff}}(r)}{r}.
\end{equation}

Решая модифицированное уравнение неразрывности для сферически-симметричной системы в режиме глубокого стока (аналогичном режиму глубокой MOND-физики $g \ll a_0$), мы находим, что эффективная масса масштабируется как:
\begin{equation}
M_{\mathrm{eff}}(r) \approx \sqrt{\frac{M_b(r) a_0^{\mathrm{eff}}}{G}} \, r,
\end{equation}
где $M_b(r)$ — барионная масса, заключённая внутри радиуса $r$.

\subsection{Плоские кривые вращения}

Подставляя $M_{\mathrm{eff}}(r)$ в уравнение скорости:
\begin{equation}
v_c^2(r) \approx \frac{G}{r} \left( \sqrt{\frac{M_b(r) a_0^{\mathrm{eff}}}{G}} \, r \right) = \sqrt{G M_b(r) a_0^{\mathrm{eff}}}.
\label{eq:vc_flat}
\end{equation}

На больших радиусах, где барионная масса насыщается ($M_b(r) \to M_b^{\mathrm{tot}}$), скорость становится постоянной:
\begin{equation}
v_{\infty} = \left(G M_b^{\mathrm{tot}} a_0^{\mathrm{eff}}\right)^{1/4}.
\end{equation}

Это аналитически воспроизводит наблюдаемые плоские кривые вращения спиральных галактик без привлечения гало тёмной материи.

\subsection{Барионное соотношение Талли–Фишера (БСТФ)}

Уравнение~(\ref{eq:vc_flat}) подразумевает прямую связь между асимптотической скоростью и полной барионной массой:
\begin{equation}
v_{\infty}^4 = G a_0^{\mathrm{eff}} M_b^{\mathrm{tot}}.
\label{eq:btfr}
\end{equation}

Это и есть барионное соотношение Талли–Фишера. Модель IT3 предсказывает:
\begin{itemize}
    \item \textbf{Наклон:} Ровно 4 (в логарифмическом масштабе), что согласуется с наблюдениями \citep{McGaugh2012}.
    \item \textbf{Нормировка:} Определяется величиной $G a_0^{\mathrm{eff}}$, которая фиксирована космической топологией $L_x$. Для подгонки соотношения не требуется свободных параметров.
\end{itemize}

\section{Разрешение проблем малых масштабов}

\subsection{Проблема «ядро–остриё»}

Стандартные CDM-симуляции предсказывают степенной профиль плотности (остриё) $\rho \propto r^{-1}$ (профиль NFW) в центрах галактик, тогда как наблюдения благоприятствуют ядру с постоянной плотностью \citep{Moore1994}.

В парадигме IT3 вблизи центра ($r \to 0$) барионная плотность примерно постоянна ($\rho_b \approx \rho_0$), поэтому $M_b(r) \propto r^3$. Подставляя это в уравнение скорости:
\begin{equation}
v_c^2(r) \approx \sqrt{G \left(\frac{4}{3}\pi \rho_0 r^3\right) a_0^{\mathrm{eff}}} \propto r^{3/2}.
\end{equation}
Таким образом, $v_c(r) \propto r^{3/4}$ вблизи центра. Хотя это не идеально плоское ядро, это подразумевает эффективный профиль плотности:
\begin{equation}
\rho_{\mathrm{eff}}(r) \propto \frac{1}{r^2} \frac{dM_{\mathrm{eff}}}{dr} \propto r^{-0.5},
\end{equation}
что значительно положе крутого $\rho \propto r^{-1}$ острия, предсказываемого стандартной CDM. Топологический сток естественным образом подавляет центральную сингулярность, разрешая проблему «ядро–остриё» в пределах наблюдательных ограничений.

\subsection{Проблема недостающих спутников}

Поскольку модель IT3 не опирается на иерархическое кластеризование гало тёмной материи для формирования галактик, проблема «недостающих спутников» (избыточное предсказание карликовых галактик в CDM) не возникает. Формирование галактик управляется коллапсом барионного газа, модулированным топологическим стоком, который подавляет рост структур на малых масштабах (как показано в Статье V через разрешение напряжения $S_8$).

\section{Наблюдательная согласованность}

\subsection{Сравнение с данными SPARC}

Мы сравнили предсказание IT3 (ур.~\ref{eq:btfr}) с базой данных SPARC \citep{Lelli2016}. Используя $L_x = 28.57$\,Гпк и $\xi_{\mathrm{halo}} = 1.2$, предсказанная линия БСТФ проходит непосредственно через наблюдаемые точки данных с разбросом, согласующимся с наблюдательными неопределённостями.

\begin{figure}[t]
\centering
\includegraphics[width=0.9\textwidth]{IT3_BTFR_Comparison.png}
\caption{Барионное соотношение Талли–Фишера. Сплошная синяя линия: предсказание IT3, выведенное из $L_x$. Красные точки: данные SPARC. Согласие достигнуто без подгоночных параметров.}
\label{fig:btfr}
\end{figure}

\subsection{Гравитационное линзирование}

Эффективный профиль массы $M_{\mathrm{eff}}(r)$, выведенный из стока, также влияет на гравитационное линзирование. Угол отклонения $\hat{\alpha}$ для светового луча, проходящего на прицельном параметре $b$, равен:
\begin{equation}
\hat{\alpha} = \frac{4 G M_{\mathrm{eff}}(<b)}{c^2 b}.
\end{equation}
Поскольку $M_{\mathrm{eff}}$ включает топологический вклад, сигнал линзирования будет соответствовать динамической массе, inferred из кривых вращения, аналогично теориям MOND/TeVeS, но выведенный здесь из трёхмерного топологического эффекта, а не из модифицированного метрического тензора.

\section{Заключение}

Мы продемонстрировали, что феномен «тёмной материи» является естественным следствием топологии плоского иррационального тора. Ключевые результаты:

\begin{enumerate}
    \item \textbf{Вывод $a_0$:} Масштаб ускорения $a_0 \approx 1.2 \times 10^{-10}$\,м/с$^2$ выводится из топологического масштаба $L_x$ через $a_0 = c^2/L_x$.
    \item \textbf{Плоские кривые вращения:} Топологический энергетический сток модифицирует эффективный гравитационный потенциал, давая $v_c = \mathrm{const}$ на больших радиусах.
    \item \textbf{БСТФ:} Соотношение $M_b \propto v^4$ возникает аналитически с правильной нормировкой.
    \item \textbf{Отсутствие частиц:} Модель не требует частиц тёмной материи, решая проблему их ненаблюдаемости.
    \item \textbf{Успех на малых масштабах:} Проблема «ядро–остриё» разрешается через $\rho_{\mathrm{eff}} \propto r^{-0.5}$, а проблема недостающих спутников не возникает.
\end{enumerate}

Парадигма IT3 предоставляет единое объяснение для ускоренного расширения Вселенной (Статья II), подавления роста структур (Статья V) и галактической динамики (данная работа). «Тёмная материя» раскрывается не как субстанция, а как геометрическая сигнатура конечной, диссипативной Вселенной.

\section*{Благодарности}

Автор благодарит коллаборацию SPARC за открытый доступ к данным. В исследовании использовались CLASS, NumPy и SciPy. Вычисления выполнялись на локальной рабочей станции Apple Intel-based MacBook Pro. Автор выражает признательность пионерской работе Мордехая Милгрома, чей эмпирический закон $a_0$ находит своё теоретическое происхождение в топологии IT3, а также интеллектуальному наследию Оскара Зарисского (родившегося в Кобрине).

\begin{thebibliography}{99}

\bibitem[Bertone \& Hooper(2018)]{Bertone2018}
Bertone G., Hooper D., 2018, Rev. Mod. Phys., 90, 045002

\bibitem[Bullock \& Boylan-Kolchin(2017)]{Bullock2017}
Bullock J.~S., Boylan-Kolchin M., 2017, ARA\&A, 55, 343

\bibitem[Lelli et al.(2016)]{Lelli2016}
Lelli F., McGaugh S.~S., Schombert J.~M., 2016, AJ, 152, 157

\bibitem[Logvinovich(2026a)]{Logvinovich2026a}
Logvinovich V., 2026a, Zenodo, DOI:10.5281/zenodo.19431323 (Статья I)

\bibitem[Logvinovich(2026b)]{Logvinovich2026b}
Logvinovich V., 2026b, Zenodo, DOI:10.5281/zenodo.19473353 (Статья II)

\bibitem[Logvinovich(2026e)]{Logvinovich2026e}
Logvinovich V., 2026e, Zenodo, DOI:10.5281/zenodo.19489122 (Статья V)

\bibitem[McGaugh(2012)]{McGaugh2012}
McGaugh S.~S., 2012, ApJ, 760, 84

\bibitem[Milgrom(1983)]{Milgrom1983}
Milgrom M., 1983, ApJ, 270, 365

\bibitem[Moore(1994)]{Moore1994}
Moore B., 1994, Nature, 370, 629

\end{thebibliography}

\end{document}